\documentclass[10pt,a4paper]{amsart}
\usepackage[margin=19mm]{geometry}
\usepackage{amsmath,amssymb,mathtools}
\usepackage[scr=rsfs,cal=euler]{mathalfa}
\usepackage{xfrac}
\usepackage[unicode,hidelinks,bookmarks=false]{hyperref}
\newcommand{\ie}{i.e.\ }
\newcommand{\cf}{cf.\ }
\newcommand{\resp}{respectively\ }
\newcommand{\Subsection}{\subsection}
\providecommand{\nobreakdash}{\nobreak\hbox{-}}
\setlength{\emergencystretch}{2em}
\multlinegap0pt
\usepackage{amssymb}
\newtheorem{theorem}{Theorem}
\title{The Maximum of the Volume of a Cevian Simplex and its Parts}
\author{Yagub N. Aliyev}
\date{Source: arXiv:2507.10067v2 (CC BY 4.0)}
\begin{document}
\maketitle

\noindent\emph{Source and licence.} The Maximum of the Volume of a Cevian Simplex and its Parts. Version arXiv:2507.10067v2 (CC BY 4.0), distributed by arXiv under the Creative Commons Attribution 4.0 International licence, http://creativecommons.org/licenses/by/4.0/, as recorded in the arXiv licence metadata for this exact version and shown on the abstract page https://arxiv.org/abs/2507.10067v2. Full licence notice: \texttt{licenses/arxiv-2507.10067-CC-BY-4.0.txt}.\par\smallskip

\noindent\textit{Editorial scope.} Counted unit: the article's theorem theorem (source0.tex lines 77--82). The article's complete original proof of it is delivered with it (source0.tex lines 83--93, one \texttt{proof} environment of 11 lines). No mathematics is written, repaired or re-derived: the statement and the proof are the article's own lines, in the article's own notation, and no retained line is substituted or re-flowed.  No further source-local context is needed: every cross-reference of every retained line resolves inside the two delivered spans.  Nothing was omitted from any retained span, so the package carries no omission record.  No retained line contains a citation, so the wrapper carries no bibliography and no bibliography entry.  No retained line contains a figure environment, an image-inclusion command or drawing code, so no image is delivered and the package declares no asset dependency.  Editorial formatting only, in addition: an AMS \texttt{amsart} preamble that restates the article's own declarations and macros used by the retained lines, namely the environments \texttt{theorem} on separate counters, together with the standard packages those lines need.  The retained environments print in this document as Theorem 1 in order of appearance, whereas the article prints the same items with the numbering of its own full document.  The complete native source of the article as distributed in the arXiv e-print for this exact version is bundled unchanged as \texttt{sources/181526/source0.tex}.
\begin{theorem} With the notation from Theorem 1,
$$[MN_1N_2\ldots N_{n}]\le\frac{(n-1)^2}{(n-\theta_n)^{n+3}}\cdot[A_1A_2\ldots A_{n+1}],\eqno(3)
$$
where $\theta_n=\frac{n+1-\sqrt{n^{2}+2 n-3}}{2}$. Equality
happens when $M$ has barycentric coordinates $(\theta_n,\theta_n,\ldots,\theta_n,1-n\theta_n)$ in the given order.
\end{theorem}

\begin{proof}
Since the function $F(\lambda_1,\lambda_2,\ldots,\lambda_{n+1})=\lambda_{n+1}\prod_{i=1}^n\frac{\lambda_i}{1-\lambda_i}$ is continuous on the compact simplex
$A_1A_2\ldots A_{n+1}$, $F$ attains its extremum values in this simplex. Note that on the boundary of the simplex
$A_1A_2\ldots A_{n+1}$ the function $F$ needs to be defined as $F = 0$ to guarantee its continuity for these points, too. As the
minimum of $F = 0$ is attained on the faces of the simplex, the maximum occurs in the interior of the simplex. By the Method of Lagrange Multipliers, for a fixed ${\lambda_i}+{\lambda_j}$ the product $\frac{\lambda_i}{1-\lambda_i}\cdot \frac{\lambda_j}{1-\lambda_j}$ reaches the maximum when $\lambda_i={\lambda_j}$. Therefore, for fixed $\sum_{i=1}^n{\lambda_i}=nx$ the product $\prod_{i=1}^n\frac{\lambda_i}{1-\lambda_i}$ reaches the maximum when $\lambda_1={\lambda_2}=\cdots=\lambda_n=x$. Then $\lambda_{n+1}=1-nx$ and $$\lambda_{n+1}\cdot\prod_{i=1}^n\frac{\lambda_i}{1-\lambda_i}=\left(\frac{x}{1-x}\right)^{n}\cdot \left(1-nx\right)=:f(x).$$
We obtain the maximum of $f$ in the interval $\left(0,\frac{1}{n}\right)$ by standard single variable calculus. The derivative $$f'(x)=\left(\frac{x}{1-x}\right)^{n} \cdot\frac{n \left(x^{2}-(n+1)x+1\right)}{x \left(1-x\right)},$$leads
to the critical point $x=\theta_n$ being the only global maximum point of $f$
in the reference interval. Therefore, $$f(x)\le f \left(\theta_n\right)=\frac{(n-1)^2}{(n-\theta_n)^{n+3}},$$
as needed. The last expression is obtained from the earlier definition
of $f$, using three equivalent forms of the quadratic equation satisfied by $\theta_n$: $1-n\theta_n=\theta_n(1-\theta_n)$, $1-\theta_n=\theta_n(n-\theta_n)$, and $(1-\theta_n)^2=(n-1)\theta_n$.
\end{proof}

\end{document}
